\begin{figure}[tbp]
  \centering
  \begingroup
  \definecolor{hierarchyblue}{RGB}{37,87,128}
  \begin{tikzpicture}[
    x=1cm,y=1cm,
    font=\small,
    coarse/.style={draw=black,line width=.8pt,rounded corners=5pt},
    fine/.style={draw=black!65,dashed,line width=.65pt,rounded corners=4pt},
    vertex/.style={circle,fill=black,inner sep=1.5pt},
    levelarrow/.style={-{Stealth[length=1.8mm]},line width=.55pt},
    every node/.style={inner sep=2pt}
  ]
    \node at (6.35,6.25) {Selected threshold partitions: fine to coarse};
    \node (pzero) at (1.10,5.78) {$\mathcal P_0$};
    \node (dotsleft) at (2.85,5.78) {$\cdots$};
    \node (pi) at (4.60,5.78) {$\mathcal P_i$};
    \node (pinext) at (7.00,5.78) {$\mathcal P_{i+1}$};
    \node (dotsright) at (9.00,5.78) {$\cdots$};
    \node (root) at (11.05,5.78) {$\{V\}$};
    \draw[levelarrow] (pzero) -- (dotsleft);
    \draw[levelarrow] (dotsleft) -- (pi);
    \draw[levelarrow] (pi) -- (pinext);
    \draw[levelarrow] (pinext) -- (dotsright);
    \draw[levelarrow] (dotsright) -- (root);

    \draw[fine] (1.20,5.12) -- (2.00,5.12);
    \node[anchor=west] at (2.10,5.12) {parts of $\mathcal P_i$};
    \draw[coarse] (7.15,5.12) -- (7.95,5.12);
    \node[anchor=west] at (8.05,5.12) {parts of $\mathcal P_{i+1}$};

    \draw[coarse] (.25,1.40) rectangle (6.15,4.65);
    \draw[coarse] (6.55,1.40) rectangle (12.45,4.65);
    \node at (3.20,4.33)
      {$C=A^*=(A')^*$};
    \node at (9.50,4.33) {$D$};

    \draw[fine] (.65,1.85) rectangle (2.95,3.98);
    \draw[fine] (3.45,1.85) rectangle (5.75,3.98);
    \draw[fine] (6.95,1.85) rectangle (9.25,3.98);
    \draw[fine] (9.75,1.85) rectangle (12.05,3.98);
    \node at (1.80,3.65) {$A$};
    \node at (4.60,3.65) {$A'$};
    \node at (8.10,3.65) {$D_1$};
    \node at (10.90,3.65) {$D_2$};

    \node[vertex] at (1.20,2.48) {};
    \node[vertex] at (1.45,3.10) {};
    \node[vertex] (u) at (2.25,2.80) {};
    \node[vertex] (v) at (4.15,2.80) {};
    \node[vertex] at (4.95,3.10) {};
    \node[vertex] (w) at (5.10,2.35) {};
    \node[vertex] (z) at (7.60,2.35) {};
    \node[vertex] at (7.70,3.10) {};
    \node[vertex] at (8.60,2.75) {};
    \node[vertex] at (10.35,2.50) {};
    \node[vertex] at (10.60,3.10) {};
    \node[vertex] at (11.40,2.80) {};

    \draw[hierarchyblue,line width=1.3pt] (u) -- (v)
      node[midway,above=3pt,text=hierarchyblue] {$e$};
    \draw[black,line width=.7pt] (w) -- (z)
      node[midway,above=3pt] {$f$};

    \node[text=hierarchyblue] at (6.35,.86)
      {$e\in B_i\setminus B_{i+1},\qquad e\in O(A)\cap O(A')$};
    \node at (6.35,.30)
      {$O(A)=E_H(A,C\setminus A),\qquad f\in B_{i+1}\subseteq B_i$};
  \end{tikzpicture}
  \endgroup
  \caption{At a selected transition $\mathcal P_i\to\mathcal P_{i+1}$,
    the parts $A,A'$ merge into their immediate parent $C$.
    The edge $e$ leaves the crossing set and belongs to both $O(A)$ and
    $O(A')$, whereas $f$ remains a crossing edge.
    Only representative vertices and edges are shown.}
  \label{fig:hierarchy}
\end{figure}
